\begin{afterword}
\summaryhead

The post, closely based on a textbook excerpt, answers critics who say that social science finds only what anyone would expect. The author calls that expectation hindsight: people who know an answer believe they would have guessed it. The post lists five findings about American soldiers in the Second World War, each with a plausible explanation, and asks how many the reader could have predicted. Then it reveals that each is the reverse of what was found, and tells a reader who noticed nothing that they have learned how good their model really was.

A study by Daphna Baratz is cited in which students rated both a finding and its opposite as what they would have predicted. Hindsight, the post concludes, makes scientific findings seem less surprising than they are, especially in fields we understand. This undervalues researchers and hides evidence that does not fit what we really expected. The remedy is a conscious effort to be shocked enough.

\responsehead

The post reports a real effect with a good demonstration. Lazarsfeld's reversed list is a fair source, and the post credits it. The Baratz study, in which students found a finding and its opposite equally predictable, is the right kind of evidence for the claim that the feeling of obviousness is unreliable.

Its claims go further than that evidence. Its reply, that the critics' ``expectation'' is ``all hindsight,'' implies that nothing social science finds could have been expected. The studies show something narrower: people cannot tell, afterwards, what they would have expected. Lazarsfeld's own account of his list is also looser than ``the opposite.'' For the climate item he reports no difference.

The lesson drawn for the reader asks more than the demonstration allows. ``The measure of your strength as a rationalist is your ability to be more confused by fiction than by reality'' repeats the thesis of ``Your Strength as a Rationalist,'' and has the problem noted there: confusion can come from a false model as easily as from a false story. Lazarsfeld's conclusion was that since every kind of human reaction is conceivable, only data can say which occurs. A reader who could not tell the reversed list from the true one has shown that they lacked data on soldiers, which was Lazarsfeld's point. Attention to their own confusion would not have supplied it.

The extension to ``the discoveries we understand,'' and with ``probably'' to neurology and physics, comes with no evidence, and the research on expertise and hindsight is mixed.

The remedy does not follow the research. The post ends by asking for ``a conscious effort to be shocked enough.'' Research on hindsight bias suggests that knowing about it and meaning to avoid it do not remove it, and that what reduces it is explaining how the other outcome could have come about. The post's own middle question, whether rationalizing either side feels different, points toward that method. Its conclusion does not take it up.

In short: a fair demonstration of a real bias, extended to the claim that ``all'' expectation is hindsight, turned into a lesson about trusting one's confusion when only data could have helped, and closed with a call for ``a conscious effort to be shocked'' where the research recommends explaining how the other outcome could have happened.
\end{afterword}
